\section{Giới thiệu}
\subsection{Động cơ}
Du lịch là một trong những ngành kinh tế phát triển nhanh chóng nhất trên thế giới. Du lịch cũng là một ngành công nghiệp kiếm có được ngoại hối mà không cần xuất khẩu của cải, hàng hóa của quốc gia \cite{vietnamtourism_export}. Trên Diễn đàn kinh tế thế giới, các sáng kiến Chuyển đổi số trong du lịch dự kiến sẽ đóng góp tới 305 tỷ USD lợi nhuận trong năm 2025. Sự phát triển của công nghệ đã cung cấp những công cụ ứng dụng và nền tảng công nghệ trong việc mang lại các trải nghiệm du lịch mới mẻ, hấp dẫn và tiện lợi với khách du lịch.

Nhận định này được đưa ra trên cơ sở hai nghiên cứu năm 2024 của booking.vn gồm: Nghiên cứu Xu hướng du lịch và Dự đoán du lịch \cite{qdnd_genz}. Theo đó, 51\% số khách du lịch thuộc thế hệ Gen Z Việt Nam tham gia khảo sát cho biết họ cảm thấy rất thoải mái khi dựa vào công nghệ nhằm tạo ra những trải nghiệm du lịch khó quên. 50\% chia sẻ rằng, họ tin tưởng AI trong việc đề xuất những địa điểm du lịch ít xô bồ và không nhiều người biết đến. 67\% sẽ sử dụng công nghệ lập kế hoạch du lịch bằng AI để lên lịch trình cho chuyến đi của mình. 57\% thể hiện sự quan tâm đến các tiện ích và dịch vụ ứng dụng công nghệ hiện đại...

61\% khách du lịch Gen Z tại Việt Nam cho biết, yếu tố chính thúc đẩy họ đi du lịch là để thư giãn. 29\% mong muốn kết nối sâu hơn với bản thân. Điều này lý giải sự lên ngôi của các loại hình du lịch tập trung vào sức khỏe, với xu hướng tìm đến những điểm đến cung cấp dịch vụ như yoga, thiền định và liệu trình trị liệu spa...

Theo Tổ chức Du lịch Liên Hợp Quốc, du lịch Việt Nam có tốc độ "tăng trưởng mạnh nhất thế giới" trong 6 tháng đầu năm 2025 \cite{vnexpress_growth}. Theo số liệu thống kê từ Cục Du lịch Quốc gia Việt Nam đón gần 11 triệu lượt, tăng 17\% so với cùng kỳ 2024 và gấp 1,3 lần so với cùng kỳ 2019. Tám tháng đầu năm, khách quốc tế đến Việt Nam đạt gần 14 triệu người, tăng gần 22\% so với cùng kỳ năm trước. Cùng với sự phát triển bùng nổ này, Bộ VHTTDL ban hành Chương trình đẩy mạnh chuyển đổi số du lịch thông minh trong lĩnh vực du lịch \cite{vietnamtourism_digital}. Chương trình hướng tới ứng dụng khoa học công nghệ trong thúc đẩy phát triển du lịch trở thành ngành kinh tế mũi nhọn, đảm bảo phát triển bền vững và hội nhập quốc tế, đóng góp vào quá trình phát triển Chính phủ số, kinh tế số và xã hội số. Mục tiêu đề ra là phát triển và hoàn thiện các nền tảng số hỗ trợ công tác quản lý nhà nước, phục vụ khách du lịch, phát triển hệ thống dữ liệu làm nền tảng cho xây dựng mô hình du lịch thông minh.

Theo thông tin từ Viện Nghiên cứu phát triển du lịch (Cục Du lịch quốc gia Việt Nam), hiện nay, du lịch giới trẻ đang bùng nổ và là một ngách đầy triển vọng phát triển của ngành du lịch \cite{qdnd_genz}. Nhiều báo cáo và dữ liệu khẳng định đây là thế hệ giàu tiềm năng du lịch, họ đi du lịch để được trải nghiệm khám phá và thử thách bản thân, bên cạnh đó là tìm kiếm sự thư giãn và kết nối cá nhân. 

Đặc biệt, xu hướng du lịch tự túc của giới trẻ và thế hệ Z ngày càng phổ biến, đi cùng nhu cầu trải nghiệm mới mẻ, cá nhân hóa và tiết kiệm chi phí. Cùng có nhận định này với đối tượng cụ thể hơn là Gen Z (những người sinh ra từ năm 1997 đến năm 2012), ông Varun Grover, Giám đốc Quốc gia của booking.com tại Việt Nam cho biết: “Sở thích du lịch của Gen Z tại Việt Nam mang đến một cái nhìn thú vị về tương lai của ngành du lịch. Tại đó, công nghệ, giá trị và sự gắn kết trong mối quan hệ có mối liên hệ mật thiết với nhau. Đối với thế hệ này, du lịch không chỉ dừng lại ở việc khám phá những vùng đất mới mà còn là cách họ thể hiện bản sắc cá nhân với sự hỗ trợ của trí tuệ nhân tạo (AI) và các công cụ kỹ thuật số giúp tối ưu hóa trải nghiệm cá nhân hóa. Dù thế hệ này đang tìm kiếm những trải nghiệm chân thực, cân bằng giữa du lịch cá nhân và kết nối gia đình hay theo đuổi những giá trị độc đáo, Gen Z đang định hình một bức tranh du lịch mang tính cá nhân hóa và sống động hơn bao giờ hết” \cite{qdnd_genz}.

Sự phát triển bùng nổ của ngành du lịch Việt Nam đi kèm với nhiều thách thức trong việc chuyển đổi số mà chưa có các biện pháp giải quyết những khó khăn, thách thức trong nhu cầu cải thiện hạ tầng và dịch vụ để thu hút và phát huy tiềm năng của tệp khách hàng này.

Thứ nhất, sự quá tải và phân mảnh thông tin khiến du khách gặp khó khăn khi lập kế hoạch. Thông tin du lịch hiện nay được phân tán trên nhiều kênh như mạng xã hội, OTA (Online Travel Agency) hay các blog du lịch, nhưng thường thiếu cấu trúc, khó xác thực và ít được cá nhân hóa.

Thứ hai, khủng hoảng niềm tin từ đánh giá giả và quảng cáo trá hình (“seeding”) đang làm xói mòn trải nghiệm của du khách. Không ít trường hợp du khách thất vọng khi thực tế khác xa so với thông tin quảng bá, gây lãng phí thời gian và chi phí.

Thứ ba, các ứng dụng hiện tại chủ yếu giải quyết từng khâu: đặt phòng (Booking, Agoda), tìm địa điểm ăn uống (Foody), hoặc truyền cảm hứng (TikTok, Instagram). Người dùng chưa có công cụ liền mạch để đi từ khâu tìm cảm hứng → lập lịch trình tối ưu → đặt dịch vụ → chia sẻ trải nghiệm.

Thứ tư, Doanh nghiệp lữ hành, dịch vụ nhà hàng, khách sạn, vận tải du lịch gặp nhiều hạn chế trong việc quảng bá dịch vụ, quản lý booking và tiếp nhận phản hồi khách hàng. Các homestay, quán cà phê, hay nhà hàng địa phương thiếu công cụ phân tích dữ liệu và xác thực thông tin, dẫn đến khó cạnh tranh với các tập đoàn lớn.

Với mục tiêu góp phần vào sự tăng trưởng và chuyển đổi số trong ngành du lịch,
Từ những vấn đề trên, xuất hiện nhu cầu cấp thiết về một hệ thống du lịch thông minh, cá nhân hóa trải nghiệm cho du khách, đồng thời hỗ trợ Doanh nghiệp lữ hành, dịch vụ nhà hàng, khách sạn, vận tải quản lý và phát triển bền vững.

\subsection{Mục tiêu}
Đề tài hướng đến xây dựng \textbf{Hệ thống Gợi ý hành trình du lịch thông minh dựa trên AI}, ứng dụng trong bối cảnh ngành du lịch Việt Nam đang bùng nổ và chuyển dịch mạnh mẽ sang mô hình du lịch thông minh. Hệ thống phát triển Recommendation Engine dựa trên ML/DL, ứng dụng NLP tiếng Việt để xử lý review và phản hồi được thiết kế để vượt qua hạn chế của các công cụ hiện tại trong việc lập kế hoạch du lịch như tương thích với sở thích cá nhân và ngân sách, hạn chế các địa điểm có đánh giá ảo. Hệ thống hứa hẹn sẽ có giao diện thân thiện, cho phép người dùng lên lịch trình, tuỳ chỉnh theo yêu cầu để từ đó đưa ra lịch trình phù hợp nhất với bản thân. Người dùng còn được đánh giá dịch vụ đã trải nghiệm và chia sẻ lịch trình, trải nghiệm của bản thân trên diễn đàn chung của hệ thống. Song song, hệ thống còn cung cấp công cụ hỗ trợ các doanh nghiệp quản lý đơn hàng, quản lý dịch vụ.

\subsection{Phạm vi}
Để đảm bảo dự án khả thi và tập trung vào việc giải quyết các mục tiêu cốt lõi trong khung thời gian và nguồn lực hạn chế, phạm vi của đề tài được xác định như sau
\subsubsection{Phạm vi chức năng}
Hệ thống tập trung vào việc xây dựng các chức năng chính sau:
\begin{itemize}
    \item \textbf{Cá nhân hóa gợi ý:} Gợi ý hành trình du lịch cá nhân dựa trên địa điểm du khách quan tâm, đề xuất các hoạt động giải trí phù hợp với sở thích, ngân sách và lịch sử tương tác của người dùng.
    
\end{itemize}

\subsubsection{Phạm vi Kỹ thuật}
Các giới hạn và lựa chọn về mặt công nghệ được xác định như sau:
\begin{itemize}
    \item \textbf{Kiến trúc hệ thống}:Dự kiến xây dựng theo kiến trúc ứng dụng kiến trúc Client-Server cơ bản để đơn giản hóa việc phát triển và triển khai trong giai đoạn đầu.
    \item \textbf{Công nghệ phát triển}:
        \begin{itemize}
            \item Frontend (Giao diện người dùng): 
            \item Backend (Xử lý logic):
            \item Cơ sở dữ liệu:
        \end{itemize}
    \item \textbf{Nền tảng triển khai}:
    \item \textbf{Bảo mật}:
\end{itemize}

\subsection{Ý nghĩa của dự án}

\subsubsection{Ý nghĩa thực tiễn}
Đề tài “Phát triển hệ thống gợi ý hành trình du lịch dựa trên AI và đánh giá cộng đồng” mang lại nhiều giá trị thiết thực đối với du khách, doanh nghiệp và toàn ngành du lịch Việt Nam. Trước hết, hệ thống giúp tiết kiệm thời gian và chi phí cho du khách thông qua việc tự động hóa quá trình tìm kiếm, lên lịch trình và đánh giá chất lượng dịch vụ, từ đó giảm thiểu rủi ro khi tiếp cận các thông tin sai lệch hoặc dịch vụ kém chất lượng. Bên cạnh đó, hệ thống góp phần thúc đẩy quá trình chuyển đổi số cho các doanh nghiệp du lịch vừa và nhỏ (Doanh nghiệp lữ hành, dịch vụ nhà hàng, khách sạn, vận tải). Doanh nghiệp lữ hành, dịch vụ nhà hàng, khách sạn, vận tải có thể chủ động tiếp cận khách hàng, quản lý hoạt động kinh doanh hiệu quả hơn và gia tăng năng lực cạnh tranh. Cuối cùng, đề tài còn có ý nghĩa đối với toàn ngành khi tạo ra một mô hình du lịch thông minh kết hợp giữa công nghệ AI và đánh giá cộng đồng. Mô hình này hướng đến một môi trường du lịch minh bạch, cạnh tranh lành mạnh, trong đó chất lượng dịch vụ và trải nghiệm thực tế trở thành yếu tố quyết định.

\subsubsection{Ý nghĩa khoa học}
Không chỉ dừng lại ở ứng dụng thực tiễn, đề tài còn mang giá trị khoa học quan trọng. Thứ nhất, đây là cơ hội để ứng dụng và tùy chỉnh các thuật toán trí tuệ nhân tạo vào một bài toán đặc thù: gợi ý hành trình du lịch trong bối cảnh văn hóa và dữ liệu Việt Nam. Các mô hình học máy như Collaborative Filtering, Content-Based Filtering hay Hybrid Models, cùng với các kỹ thuật học sâu (Deep Learning), có thể được kiểm nghiệm và so sánh để lựa chọn giải pháp tối ưu. Thứ hai, đề tài mở ra hướng nghiên cứu về xử lý ngôn ngữ tự nhiên (NLP) tiếng Việt thông qua việc phân tích đánh giá, bình luận và phản hồi của du khách. Việc trích xuất cảm xúc (Sentiment Analysis) và phát hiện đánh giá giả trong ngữ cảnh tiếng Việt là một thách thức, đồng thời là đóng góp khoa học có giá trị. Thứ ba, hệ thống cho phép xây dựng một bộ dữ liệu về du lịch Việt Nam bao gồm sở thích, hành trình và đánh giá của du khách. Bộ dữ liệu này không chỉ phục vụ cho việc cải tiến hệ thống mà còn có thể trở thành nguồn dữ liệu quý cho các nghiên cứu chuyên sâu về hành vi người tiêu dùng và du lịch thông minh trong tương lai.

\subsection{Cấu trúc báo cáo}
Để trình bày một cách hệ thống và khoa học các kết quả nghiên cứu và phát triển, báo cáo đồ án được tổ chức thành 8 chương chính:

\textbf{Chương 1 - Giới thiệu}: Đặt vấn đề, trình bày bối cảnh, lý do chọn đề tài, các thách thức gặp phải, từ đó xác định mục tiêu, phạm vi cụ thể, và ý nghĩa của đồ án.

\textbf{Chương 2 - Cơ sở lý thuyết}: Cung cấp nền tảng kiến thức về nghiệp vụ du lịch và các công nghệ, kỹ thuật được sử dụng để phát triển hệ thống.

\textbf{Chương 3 - Các hệ thống liên quan}: Khảo sát, phân tích và so sánh các giải pháp gợi ý hành trình du lịch hiện có trên thị trường, nhằm rút ra các bài học kinh nghiệm và xác định điểm mới của hệ thống.

\textbf{Chương 4 - Tổng quan hệ thống}: Đi sâu vào phân tích và đặc tả các yêu cầu của hệ thống, bao gồm bối cảnh nghiệp vụ, yêu cầu người dùng, và các luồng chức năng chính.

\textbf{Chương 5 - Thiết kế hệ thống}: Trình bày chi tiết các quyết định thiết kế, bao gồm kiến trúc hệ thống, thiết kế cơ sở dữ liệu, sitemap, và thiết kế giao diện người dùng (UI/UX).

\textbf{Chương 6 - Hiện thực \& Kiểm thử}: Mô tả quá trình hiện thực hóa các thiết kế, cấu trúc mã nguồn, và trình bày quy trình kiểm thử nhằm đảm bảo chất lượng phần mềm.

\textbf{Chương 7 - Đánh giá}: Đánh giá toàn diện các kết quả đạt được, đối chiếu với các yêu cầu chức năng và phi chức năng đã đề ra ban đầu.

\textbf{Chương 8 – Tổng kết}: Tóm tắt lại toàn bộ quá trình thực hiện đồ án, nêu bật những đóng góp chính, đồng thời đề xuất những hướng phát triển trong tương lai.

% [1]https://vietnamtourism.gov.vn/post/42949
% [2]https://vnexpress.net/du-lich-viet-nam-tang-truong-manh-nhat-the-gioi-4938099.html
% [3]https://vietnamtourism.gov.vn/post/63900
% [4]https://www.qdnd.vn/du-lich/cac-van-de/nhin-nhan-tuong-lai-du-lich-viet-qua-so-thich-cua-gioi-tre-802302
% [5]https://www.qdnd.vn/du-lich/cac-van-de/nhin-nhan-tuong-lai-du-lich-viet-qua-so-thich-cua-gioi-tre-802302
